% !TEX program = pdflatex
  % !BIB program = bibtex

\documentclass[sigconf]{acmart}
\usepackage{balance}

\title{An Evaluation of The Performance of \\[0.5em]
Lightweight Language Models on NL2SQL tasks}

\author{Tom Fan}
\affiliation{
  \institution{Carleton University}
  \city{Ottawa}
  \country{Canada}
}
\email{tomfan3@cmail.carleton.ca}

\begin{document}

\maketitle


\section{Abstract}

Since their introduction as a consumer product in 2022, transformer-based large language models have been used for 
automating many kinds of digital tasks. This includes natural language to code and SQL tasks, which have exploded in 
research and development in the last few
years. While a large proportion of LLM use is likely done using proprietary, closed source, and large footprint 
foundation models hosted by SAAS providers like Anthropic or OpenAI, an increasing amount of open-source models are also 
entering the market. These models have demonstrated promise on several different kinds of tasks, such as code generation
using Python \cite{hasan2025smallcode}. In this paper, we aim to provide an overview of how several of the smaller open
source models perform when faced with elementary NL2SQL tasks on the Spider-Dev question dataset.

\section{Introduction}

In recent years, there has been a significant amount of high performing open source LLMs available from both research 
institutions and tech companies \cite{llmstats2025}. While many of the top-ranking models require expensive and complex
hardware setups to run as their parameter counts run in the hundreds of billions or even trillions, many are either 
small footprint by design or distilled versions of larger source models. 

In these cases, consumer level hardware should be able to run the models using tools like localLLama, enabling their use
offline and without dependency on external SAAS providers (for purposes of this paper, we define consumer level hardware
to be current generation Nvidia RTX chips or other equivalents, able to run models in the 20-30b parameter range or smaller
at 8 or 16b quantization). 

At the same time, there has been significant development in the field of NL2SQL tools. While previous generations of 
NL2SQL tools depended on hardcoded rules or deep neural networks \cite{nl2sqlsurvey}, the emergence of LLMs has enabled
a whole new paradigm of NL2SQL tools based on them. In fact, in their MVP state, a LLM itself can serve as a NL2SQL tool
on its own. 

\section{Testing Methodology}

\section{Results}

\section{Analysis}

\section{Conclusion}


\balance
\bibliographystyle{ACM-Reference-Format}
\bibliography{references}

\end{document}